%%%%%%%%%%%%%%%%%%%%%%%%%%%%%%%%%%%%%%%%%%%%%%%%%%%%%%%%%%%%%%%%%%%%%%%%%%%
% Función: portada interior blanca de TFG/TFM de la ETSIT-UPM.
%
% Uso: inclusión automática; normalmente no debes modificarlo.
% Modelo: 01-Plantilla-2-conAnexos-2_MODIFICADO_GRADOS-2.docx.
% No se compila por separado.
%%%%%%%%%%%%%%%%%%%%%%%%%%%%%%%%%%%%%%%%%%%%%%%%%%%%%%%%%%%%%%%%%%%%%%%%%%%

\thispagestyle{empty}
\begingroup
\providecommand{\UPMETsitCoverDegree}{\myDegreefullwrappedUpcase}
\providecommand{\UPMETsitCoverWorkType}{\MakeUppercase{\myWorkTypeFull}}
\begin{tikzpicture}[remember picture,overlay]
  \node[anchor=north,align=center,text width=18cm,inner sep=0pt,
    font=\sffamily\bfseries\fontsize{20}{24}\selectfont]
    at ([xshift=10.5cm,yshift=-2.5cm]current page.north west)
    {\MakeUppercase{\myUniversity}};
  \node[anchor=north,align=center,text width=18cm,inner sep=0pt,
    font=\sffamily\bfseries\fontsize{16}{19}\selectfont]
    at ([xshift=10.5cm,yshift=-4cm]current page.north west)
    {ESCUELA TÉCNICA SUPERIOR\\DE INGENIEROS DE TELECOMUNICACIÓN};
  \node[anchor=north,inner sep=0pt]
    at ([xshift=10.5cm,yshift=-6cm]current page.north west)
    {% Los JPG originales son cuadrados y contienen amplios márgenes blancos.
     \adjustbox{trim=0 {0.4\height} 0 {0.32\height},clip,width=8cm}{%
       \ifthenelse{\equal{\colorspaceused}{rgb}}
         {\includegraphics{logos/upm/logoetsit_tamano.jpg}}
         {\includegraphics{logos/upm/logoetsit_tamano-cmyk.jpg}}}};
  \node[anchor=north,align=center,text width=16cm,inner sep=0pt,
    font=\sffamily\bfseries\fontsize{20}{24}\selectfont]
    (upm-work-type) at ([xshift=10.5cm,yshift=-10cm]current page.north west)
    {\UPMETsitCoverDegree\\[0.35cm]\UPMETsitCoverWorkType};
  % El borde inferior del año queda a 2,5 cm del borde de la página.
  \node[anchor=south,align=center,text width=16cm,inner sep=0pt,outer sep=0pt,
    font=\sffamily\bfseries\fontsize{20}{24}\selectfont]
    (upm-author-year) at ([yshift=2.5cm]current page.south)
    {\myAuthorFullName\\[0.35cm]\myThesisDefenseYear};
  % Centro del título al 35 % del hueco real: más próximo al tipo de trabajo.
  \path (upm-work-type.south) -- (upm-author-year.north)
    node[pos=0.35,anchor=center,align=center,text width=15cm,inner sep=0pt,
      font=\sffamily\bfseries\fontsize{20}{24}\selectfont]
      {\hyphenchar\font=-1\hyphenpenalty=10000\sloppy\MakeUppercase{\myBookTitle}};
\end{tikzpicture}
\null
\clearemptydoublepage
\endgroup
